\begin{section}{Finite Sets} \label{sec:FiniteSets}

Let $J_{n}$ be a set defined by
\[
J_{n} \coloneqq \set*{1, 2, \dots, n} \quad \text{for each } n \in \N.
\]
A set $X$ is said to be \textbf{finite} if $J_{n} \hyperref[thm:BijectionEquivalenceRelation]{\sim} X$ for some
$n \in \N$. In this case, $n$ is called the \textbf{cardinality} of $X$, denoted by $\abs{X}$.
Especially, if $X = \emptyset$, we define $\abs{X} = 0$. $X$ is said to be \textbf{infinite} if $X$ is not finite.

\bigskip

\begin{lemma} \label{lem:SubsetOfJnIsFinite}
    Any subset of $J_{n}$ is finite for each $n \in \N$.
\end{lemma}

\begin{proof}
    Let $P(n)$ be the statement that any subset of $J_{n}$ is finite.
    \begin{itemize}
        \item \textbf{Base case:} $J_{1} = \set*{1}$ has $\emptyset$ and $J_{1}$ as its only subsets. We have
        $J_{1} \hyperref[thm:BijectionEquivalenceRelation]{\sim} J_{1}$ by $f \equiv id_{J_{1}}$, and $\emptyset$ is
        finite by definition. Thus $P(1)$ is true.
        \item \textbf{Inductive step:} Assume that $P(n)$ is true for some $n \in \N$.
        Let $A \subseteq J_{n + 1}$ be given. We have two cases:
        \begin{itemize}
            \item If $n + 1 \notin A$, then $A \subseteq J_{n}$, which is finite by the inductive hypothesis.
            \item If $n + 1 \in A$, then $A = B \cup \set*{n + 1}$ for some $B \subseteq J_{n}$.
            By the inductive hypothesis, $B$ is finite, so there exists a bijection $f: J_{m} \to B$
            for some $m \in \N$. Let $g: J_{m + 1} \to A$ be defined by
            \[
            g(k) = \begin{cases}
                f(k) & \text{if } k \leq m \\
                n + 1 & \text{if } k = m + 1
            \end{cases}
            \]
            Then $g$ is bijective, which implies that $A$ is finite.
        \end{itemize}
        Thus $P(n + 1)$ is true whenever $P(n)$ is true.
    \end{itemize}
    By the principle of mathematical induction, $P(n)$ is true for all $n \in \N$.
\end{proof}

\bigskip

\begin{theorem} \label{thm:SubsetOfFiniteSetIsFinite}
    Any subset of a finite set is finite.
\end{theorem}

\begin{proof}
    Let $X$ be a finite set and $Y \subseteq X$. Then there exists a bijection
    $f: J_{n} \to X$ for some $n \in \N$. Let $Z$ be defined by
    \[
    Z \coloneqq f^{-1}(Y) \subseteq J_{n}.
    \]
    Since $Z \subseteq J_{n}$, $Z$ is finite by
    \Cref{lem:SubsetOfJnIsFinite}, so there exists a bijection $g: J_{m} \to Z$ for some $m \in \N$.
    \\[6mm]
    By \Cref{thm:PropertiesOfRestrictionFunctions}, $f|_{Z}^{Y}$ is bijective. Then, by
    \Cref{thm:PropertiesOfCompositionFunctions}, the composition $f|_{Z}^{Y} \circ g: J_{m} \to Y$
    is also bijective. Therefore, $Y$ is finite.
\end{proof}

\bigskip

\begin{lemma} \label{lem:SurjectionFromJnIsFinite}
    Let $J_{n}$ be given for some $n \in \N$. If there exists an surjection $f: J_{n} \to Y$, then $Y$ is finite.
\end{lemma}

\begin{proof}
    Let $f: J_{n} \to Y$ be a surjective function. For any $y \in Y$, let $N_{y}$ be a set defined by
    \[
    N_{y} \coloneqq f^{-1}(\set*{y}) \subseteq J_{n}.
    \]
    Since $f$ is surjective, $N_{y}$ is nonempty. By \Cref{thm:WellOrderingProperty}, $N_{y}$ has a least element, say
    $n_{y}$. Let $M_{Y}$ be a set defined by
    \[
    M_{Y} \coloneqq \setbuilder*{n_{y}}{y \in Y} \subseteq J_{n}.
    \]
    Let $g: M_{Y} \to Y$ be a function defined by
    \[
    g(n_{y}) = y \quad \text{for each } n_{y} \in M_{Y}.
    \]
    Then $g$ is bijective. By \Cref{lem:SubsetOfJnIsFinite}, $M_{Y}$ is finite, so there exists
    a bijection $h: J_{m} \to M_{Y}$ for some $m \in \N$. Then, by \Cref{thm:PropertiesOfCompositionFunctions},
    the composition $g \circ h: J_{m} \to Y$ is bijective. Therefore, $Y$ is finite.
\end{proof}

\bigskip

\begin{theorem} \label{thm:SurjectionFromFiniteSetIsFinite}
    Let $X$ be a finite set. If there exists an surjection $f: X \to Y$, then $Y$ is finite.
\end{theorem}

\begin{proof}
    Let $f: X \to Y$ be a surjective function. Since $X$ is finite, there exists a bijection $g: J_{n} \to X$
    for some $n \in \N$. Then, by \Cref{thm:PropertiesOfCompositionFunctions}, the composition $f \circ g: J_{n} \to Y$
    is surjective. Therefore, by \Cref{lem:SurjectionFromJnIsFinite}, $Y$ is finite.
\end{proof}

\bigskip

\begin{theorem} \label{thm:FiniteOrderedSetHasExtremum}
    Let $X$ be a finite ordered set. Then $X$ has both a least element and a greatest element.
\end{theorem}

\begin{proof}
    Since $X$ is finite, there exists a bijection $f: J_{n} \to X$ for some $n \in \N$. The elements
    of $X$ can be listed as $X = \set*{a_{1}, a_{2}, \dots, a_{n}}$, where $a_{k} = f(k)$ for each $k = 1, 2, \dots, n$.
    \\[6mm]
    First, we will show that $X$ has a least element. Let $P(n)$ be the statement that
    any finite ordered set of cardinality $n$ has a least element.
    \begin{itemize}
        \item \textbf{Base case:} $X = \set*{a_{1}}$ has a least element $a_{1}$. Thus $P(1)$ is true.
        \item \textbf{Inductive step:} Assume that $P(n)$ is true for some $n \in \N$.
        Let $X$ be a finite ordered set of cardinality $n + 1$ such that
        \[
        X = \set*{a_{1}, a_{2}, \dots, a_{n + 1}}.
        \]
        By the inductive hypothesis, a finite ordered subset $Y \subseteq X$ of cardinality $n$ such that 
        \[
        Y = \set*{a_{1}, a_{2}, \dots, a_{n}}
        \]
        has a least element, say $a_{k}$ for some $k = 1, 2, \dots, n$. We have two cases:
        \begin{itemize}
            \item If $a_{k} \leq a_{n + 1}$, then $a_{k}$ is the least element of $X$.
            \item If $a_{n + 1} < a_{k}$, then $a_{n + 1}$ is the least element of $X$.
        \end{itemize}
        Thus $P(n + 1)$ is true whenever $P(n)$ is true.
    \end{itemize}
    By the principle of mathematical induction,
    $P(n)$ is true for all $n \in \N$.
    \\[6mm]
    Next, we will show that $X$ has a greatest element. Let $Q(n)$ be the statement that
    any finite ordered set of cardinality $n$ has a greatest element.
    \begin{itemize}
        \item \textbf{Base case:} $X = \set*{a_{1}}$ has a greatest element $a_{1}$. Thus $Q(1)$ is true.
        \item \textbf{Inductive step:} Assume that $Q(n)$ is true for some $n \in \N$.
        Let $X$ be a finite ordered set of cardinality $n + 1$ such that
        \[
        X = \set*{a_{1}, a_{2}, \dots, a_{n + 1}}.
        \]
        By the inductive hypothesis, a finite ordered subset $Y \subseteq X$ of cardinality $n$ such that
        \[
        Y = \set*{a_{1}, a_{2}, \dots, a_{n}}
        \]
        has a greatest element, say $a_{k}$ for some $k = 1, 2, \dots, n$. We have two cases:
        \begin{itemize}
            \item If $a_{k} \geq a_{n + 1}$, then $a_{k}$ is the greatest element of $X$.
            \item If $a_{n + 1} > a_{k}$, then $a_{n + 1}$ is the greatest element of $X$.
        \end{itemize}
        Thus $Q(n + 1)$ is true whenever $Q(n)$ is true.
    \end{itemize}
    By the principle of mathematical induction,
    $Q(n)$ is true for all $n \in \N$.
\end{proof}

\bigskip

\begin{lemma} \label{lem:PropertiesOfJmJn}
    Let $J_{m}, J_{n}$ be given for some $m, n \in \N$. Then we have the following properties:
    \begin{enumerate}[label=\textbf{(\Alph*)}]
        \item If there exists an injection $f: J_{m} \to J_{n}$, then we have $m \leq n$.
        \item If there exists a surjection $f: J_{m} \to J_{n}$, then we have $m \geq n$.
        \item If there exists a bijection $f: J_{m} \to J_{n}$, then we have $m = n$.
    \end{enumerate}
    This lemma guarantees the uniqueness of the cardinality. If $J_{m} \hyperref[thm:BijectionEquivalenceRelation]{\sim} X$
    and $J_{n} \hyperref[thm:BijectionEquivalenceRelation]{\sim} X$, then we have
    $J_{m} \hyperref[thm:BijectionEquivalenceRelation]{\sim} J_{n}$, which implies that $m = n$ by \textbf{(C)}.
\end{lemma}

\begin{proof}
    We will prove \textbf{(A)} by mathematical induction on $m$, \textbf{(B)} by mathematical induction on $n$,
    and \textbf{(C)} follows from \textbf{(A)} and \textbf{(B)}.
    \begin{enumerate}[label=\textbf{(\Alph*)}]
        \item Let $n \in \N$ be fixed.
        Let $P(m)$ be the statement that if there exists an injective function
        $f_{m}: J_{m} \to J_{n}$, then we have $m \leq n$.
        \begin{itemize} 
            \item \textbf{Base case:} Let $m = 1$. Then $m \leq n$ holds. Thus $P(1)$ is true.
            \item \textbf{Inductive step:} Assume that $P(m)$ is true for some $m \in \N$.
            If $n = 1$, there does not exist an injective function $f_{m + 1}: J_{m + 1} \to J_{n}$
            since $f_{m+1}(1) \neq f_{m+1}(m + 1)$, where $P(m + 1)$ is true.
            \\[6mm]
            Otherwise, assume that such a function $f_{m+1}$ exists. We have two cases:
            \begin{itemize}
                \item If $f_{m+1}(m + 1) = n$, let $f_{m}$ be a function defined by
                \[
                f_{m} \equiv f_{m+1}|_{J_{m}}^{J_{n - 1}}: J_{m} \to J_{n - 1}.
                \]
                Then $f_{m}$ is injective by \Cref{thm:PropertiesOfRestrictionFunctions}.
                By the inductive hypothesis, we have $m \leq n - 1$, which implies that $m + 1 \leq n$.
                \item If $f_{m+1}(m + 1) = x \neq n$, let $g_{n}: J_{n} \to J_{n}$ be defined by
                \[
                g_{n}(y) = \begin{cases}
                    x, & \text{if } y = n \\
                    n, & \text{if } y = x \\
                    y, & \text{otherwise}
                \end{cases}
                \]
                Then $g_{n}$ is bijective. Let $f_{m}$ be a function defined by
                \[
                f_{m} \equiv (g_{n} \circ f_{m+1})|_{J_{m}}^{J_{n - 1}}: J_{m} \to J_{n - 1}.
                \]
                Then $f_{m}$ is injective by \Cref{thm:PropertiesOfCompositionFunctions,thm:PropertiesOfRestrictionFunctions}.
                By the inductive hypothesis, we have $m \leq n - 1$, which implies that $m + 1 \leq n$.
            \end{itemize}
            Therefore, $P(m + 1)$ is true whenever $P(m)$ is true.
        \end{itemize}
        By the principle of mathematical induction,
        $P(m)$ is true for all $m \in \N$.
        \item Let $m \in \N$ be fixed.
        Let $Q(n)$ be the statement that if there exists a surjective function
        $f_{n}: J_{m} \to J_{n}$, then we have $m \geq n$.
        \begin{itemize}
            \item \textbf{Base case:} Let $n = 1$. Then $m \geq n$ holds. Thus $Q(1)$ is true.
            \item \textbf{Inductive step:} Assume that $Q(n)$ is true for some $n \in \N$.
            If $m = 1$, there does not exist a surjective function $f_{n + 1}: J_{m} \to J_{n + 1}$
            since $f_{n+1}(1)$ is unique, where $Q(n + 1)$ is true.
            \\[6mm]
            Otherwise, assume that such a function $f_{n+1}$ exists. We have two cases:
            \begin{itemize}
                \item If $f_{n+1}(m) = n + 1$, let $K_{n + 1}$ be a set defined by
                \[
                K_{n + 1} \coloneqq (f_{n + 1})^{-1}(\set*{n + 1}) \subseteq J_{m}.
                \]
                Let $f_{n}: J_{m} \to J_{n}$ be a function defined by
                \[
                f_{n}(y) = \begin{cases}
                    f_{n+1}(y), & \text{if } y \in J_{m} \setminus K_{n + 1} \\
                    n, & \text{if } y \in K_{n + 1}
                \end{cases}
                \]
                Then $f_{n}$ is surjective. By the inductive hypothesis, we have $m - 1 \geq n$,
                which implies that $m \geq n + 1$.
                \item If $f_{n+1}(m) = x \neq n + 1$, let $g_{n+1}: J_{n + 1} \to J_{n + 1}$ be a function defined by
                \[
                g_{n+1}(y) = \begin{cases}
                    x, & \text{if } y = n + 1 \\
                    n + 1, & \text{if } y = x \\
                    y, & \text{otherwise}
                \end{cases}
                \]